\documentclass[11pt,a4paper]{article}
\usepackage{fontspec}
\usepackage{xeCJK}
\setmainfont{msyh.ttc}[Path=C:/Windows/Fonts/,FontIndex=0,BoldFont=msyhbd.ttc]
\setCJKmainfont{msyh.ttc}[Path=C:/Windows/Fonts/,FontIndex=0,BoldFont=msyhbd.ttc]
\newfontfamily\titlelatin{simhei.ttf}[Path=C:/Windows/Fonts/,AutoFakeBold=2]
\newCJKfontfamily\titlecjk{simhei.ttf}[Path=C:/Windows/Fonts/,AutoFakeBold=2]
\usepackage[margin=22mm]{geometry}
\usepackage{graphicx}
\usepackage{needspace}
\usepackage{longtable}
\usepackage{array}
\usepackage[unicode,hidelinks,breaklinks]{hyperref}
\usepackage{fancyhdr}
\pagestyle{fancy}\fancyhf{}
\fancyfoot[C]{SRP 固定引导与原生自适应\quad\thepage}
\renewcommand{\headrulewidth}{0pt}
\setlength{\parindent}{0pt}
\setlength{\parskip}{7pt}
\linespread{1.28}
\emergencystretch=3em
\raggedbottom
\newcommand{\reviewtitle}[1]{\begin{center}{\titlelatin\titlecjk\bfseries\fontsize{21}{28}\selectfont #1}\end{center}\vspace{8pt}}
\newcommand{\reviewsection}[1]{\par\Needspace{9\baselineskip}\vspace{10pt}\begin{center}{\titlelatin\titlecjk\bfseries\fontsize{14}{20}\selectfont #1}\end{center}}
\newcommand{\reviewsubsection}[1]{\par\Needspace{4\baselineskip}\vspace{4pt}\begin{center}{\titlelatin\titlecjk\bfseries\fontsize{12}{18}\selectfont #1}\end{center}}
\begin{document}

\reviewtitle{SRP项目基石 · 设计主稿}
设计候选 SRP-\allowbreak{}FIXED-\allowbreak{}VS-\allowbreak{}ADAPTIVE-\allowbreak{}NATIVE-\allowbreak{}20261008｜用户澄清后的重构｜设计完成不等于协议已冻结。


\reviewsection{一、项目要解决的问题}
\reviewsubsection{1.1 核心定位}
面向校园学生短时情绪调节与呼吸练习，比较两种完整方案：A为固定呼吸引导圈与预设场景变化；B为依据实际数据、被试状态和场景状态进行受限调整的原生场景引导。两组追求相同理想呼吸结构与节奏；只有B形成数据返回引导与场景的闭环，并允许实际提示稍作等待或放慢，再回到理想目标。


核心问题是：这种原生自适应完整方案，能否带来值得采用的情绪体验与练习收益，其代价是什么，何时会失效？目标/\allowbreak{}实际区分、质量提示和理解检查服务于可信实现及解释结果，不再替代整项研究问题。


本轮保留校园方向、Polar H10与PLUX respiBAN BLE、Unity 2022.3.62f3及现有天气资产，只重构设计。5人、半个月、100+潜在参与者用于后续可行性分析。原有v1.2协议、API、任务状态、签收和领取快照保持原身份；本候选需要正式修订才能进入研究执行。


\reviewsubsection{1.2 产品需求与替代办法}
候选场景为学生在校园学习间隙，由操作员协助完成短时练习。任务是理解引导、舒适跟随、在暂时跟不上时获得适当支持，并能够停止。自然休息、音频引导及固定视觉节拍器均是现实替代办法。我们选择固定圈加预设场景作为可信的主动比较方案，它本身应完整、可用且具有相当的视觉品质。


待验证需求是：固定节奏可能难以跟随或显得单调；根据自身动作发生的场景响应可能提升参与感、支持继续练习，也可能带来监控压力、追逐画面、分心与额外操作成本。上述需求尚无本项目用户证据，不预设原生自适应更好。


产品采用价值还取决于穿戴、准备、校准、清洁、故障处理和维护时间。试验中两组同样佩戴设备以便公平测量；这不能估计实际部署时B相对于无需设备的A所增加的全部负担，后续实际使用评估须单独核算。


\reviewsubsection{1.3 社会价值与研究条件的边界}
价值假设是帮助校园学生完成自愿、可理解的短时练习，仍需以实际需求和使用成本验证。研究两组同佩戴设备便于比较，但现实固定引导可以无设备使用；本研究不直接证明两设备值得购置或部署。准备/\allowbreak{}校准时间、求助、拒绝和退出、管理人时及再次自主使用应在产品验证中记录。H10作为辅助研究测量的价值不自动等于产品必须配备两设备；本轮保留硬件约束，后续产品取舍另论证。

\reviewsection{二、先确定因果对象，再确定评价}
\reviewsubsection{2.1 两组究竟不同在哪里}
A：共同理想节奏I→固定引导圈和预设场景→实际呼吸X。设备记录X仅用于研究测量，不改变圈、场景、声音、奖励或模块进度。场景变化由预定时间线及固定环境随机种子控制，独立于该参与者的当前和先前生理数据。


B：共同理想节奏I、独立实际呼吸X、输入质量Q与场景状态Z→冻结的受限自适应规则π→实际下发引导G及原生场景反馈V→后续实际呼吸。B可以调整G而不能修改I来宣布用户已经达标。状态只能使用决策发生时已经获得的数据，不读取未来或事后结果。


因此处理差异同时包含引导载体、真实数据反馈及受限时序适应。两臂随机比较估计这一整包方案的分配效应；不能拆出“仅原生形式”“仅真实反馈”或“仅自适应”的独立效应。若未来需要分解机制，另设计固定原生、数据反馈但固定节奏等条件并增加资源，当前不将它们混入两组。


\reviewsubsection{2.2 可以保持相同与不应强行相同的内容}
保持相同：理想结构及初始节奏、天气资产与镜头、模块顺序分配方法、预定模块时长、设备佩戴及采集、静态环境音、教学时间预算、量表、帮助/\allowbreak{}退出规则与操作员关注。两组在处理前接受相同中性准备，教学内容须各自准确，不暗示B更先进或更有效。


有意不同：A的开环固定圈与时间驱动环境；B的原生引导、实际响应、状态相关累计变化及受限时序调整。不能再要求两组相同累计输出、相同Actual信息或每一帧相同视觉轨迹。B改变G后完成周期数可能不同，这是处理的一部分；通过固定暴露预算及实际剂量报告解释，而不是补时直到B完成同样次数。


可比性控制的是基础体验、机会和观察方式；不消除研究本来要比较的闭环差异。复杂度、显著性、位置和教学内容仍属于整包解释边界，应记录材料，不用处理后美感或理解分数调整主效应。


\reviewsubsection{2.3 三个研究问题}
RQ1：在相同理想目标及暴露预算下，B相对A产生怎样的即时情绪结果？主要效益候选保持前测校正的PANAS负性情绪后测。该选择对应产品用途并保留已有测量基础，须先验证版本、时间指令与地板分布；不是把旧条件研究直接视为新研究已冻结。


RQ2：这些情绪结果伴随怎样的理想目标达成、实际提示跟随、努力与舒适代价？把“用户跟上移动中的提示”和“用户接近共同理想节奏”分开评价。研究既允许B减少负担且接近理想，也允许B仅使任务变容易而远离目标。


RQ3：哪些状态、调整及用户解释伴随有益或有害体验，从而形成有限设计边界？用已发生的调整、信号质量、错误和访谈解释完整方案；适应次数与结果的相关不证明某次适应有因果效益。独立新实例应用属于后续规则验证，不自动成为核心研究必有的成果。


\reviewsection{三、理想目标、下发引导与实际呼吸}
\reviewsubsection{3.1 三条轨迹必须分别保留}
I是共同理想参照：模块内应练习的步骤序列、理想持续时间及相对关系，结果前固定且两组一致。G是参与者实际收到的提示及其发生时间：A按预设时钟输出，B由π受限调整。X是设备形成并经独立方法核查的实际呼吸步骤、时长与质量。不能用G补X，也不能用B调整后的G覆盖I。


现行系统只有既定Target/\allowbreak{}Actual语义，当前Target不能被悄悄解释为新的自适应目标。后续P-\allowbreak{}01/\allowbreak{}F-\allowbreak{}05/\allowbreak{}U-\allowbreak{}02/\allowbreak{}S-\allowbreak{}02等须设计I/\allowbreak{}G/\allowbreak{}X的映射、事件记录及回放方式；本轮只提出语义要求，不改消息Schema或运行API。


\reviewsubsection{3.2 被试状态与场景状态的最小定义}
首版被试状态包括可可靠识别的呼吸步骤与节奏、相对G的偏差趋势、输入质量、运动干扰，以及明确的停止/\allowbreak{}不适表达。跟不上不自动等于紧张，HRV变化不自动等于情绪变化。H10用于辅助过程观察；若要让其指标参与控制，须另有时间窗、可识别性和增量作用论证，不能先把两个设备都塞入一个情绪分数。


场景状态包括当前模块、目标/\allowbreak{}实际载体所处阶段、反馈强度是否饱和、模块剩余时间、过渡/\allowbreak{}暂停状态。它用来限制合法动作和维持可理解反馈，不能作为“画面已经变好，所以参与者已经改善”的证据。不得仅因场景尚未完成视觉变化而延长本人的练习时间。


\reviewsubsection{3.3 首版受限适应策略}
首选简单且可解释的规则控制，全部规则在结果前固定，个体参数按同一预设校准程序产生。默认先在合法周期边界决定下一周期的吸/\allowbreak{}呼时长调整及反馈强度；必要的等待表现为明确的自然呼吸过渡窗口。不得在当前吸后/\allowbreak{}呼后停留步骤中无限等待，不要求用户为补齐动画追加吸气或屏息。


当质量良好且偏差持续可确认时，在预设上下界及最大单次变化内调整下一段合法提示，使其暂时接近实际可跟随节奏。若用户重新稳定跟随，按限定回归速率使G靠近I；若仍跟不上，达到调整/\allowbreak{}缓冲上限后进入预设休息、固定指导或停止路径，不继续把目标移走。小幅波动不触发来回切换。


当输入低质量或断流时，停止个体化决策与实际成功提示；按预先规定的中性固定指导或自然呼吸路径运行。恢复后需要可用数据与合法边界才恢复适应，丢失段不补判成功。整个行为属于B的分配方案，发生降级者仍保留在B主分析集。回归理想是一项策略目标，不保证任意被试都能收敛；未回归本身属于结果。


允许调整的具体步骤、范围、更新频率、稳定观察窗、回归速度、缓冲总预算和停止条件均待技术与舒适性预试校准。本稿不填任意秒数或以现有天气秒数作为生理最优值。两组模块暴露预算固定，缓冲时间在B预算内占用，并作为适应成本报告。


\reviewsubsection{3.4 为什么要同时报告控制代价}
B可能通过降低要求获得更高的G跟随率。必须同步报告G相对I的偏移、适应持续时间、回归情况、自然呼吸缓冲占时和未完成结构。若X跟随G很好而X长期偏离I，应解释为提示迁就或任务放宽，不能称理想训练更有效。


对I的评价侧重结构、持续时间、周期节奏和固定练习预算中的达成覆盖，不用与全局理想时钟同相位作为唯一成功标准。否则一次合法缓冲就会永久记为错相。对G的同步跟随另按实际提示时间线计算；两个问题各有参照，不能混成一个PF。


\reviewsubsection{3.5 共同理想I是任务参照，不是普适最优}
I首先定义研究中的共同练习结构与节奏，不能因沿用工程预设就称其对所有人最有益。正式前逐项检查结构的用户用途、文献支持范围、舒适性与可辨识性；需要修订时在分配及正式比较前共同冻结，不能只改B或用正式结果回改。相对I更准确只证明任务达成更好，不自动证明情绪或其他收益更大；达成下降也不自动意味着总体价值更低。若多步骤无必要性，应重审范围，而不是为创新制造复杂度。

\reviewsection{四、场景与工程资产如何服务核心}
\reviewsubsection{4.1 固定组的完整体验}
保留溪边木屋Storm的雨幕、闪电、波纹与固定目标圈；圈按预设3—3—3—3等候选结构循环，其他天气按共同I。固定组关闭显示真实Actual的内圈，以及任何由本人数据驱动的累计、颜色、环境或音轨变化。现有内圈接口作为工程资产保留，不代表A必须使用双圈；禁止让内圈复制目标而被误以为真实表现。


A的背景保留有品质的预设环境变化，不通过删光美术、降低流畅度或制造无聊来衬托B。时间驱动的褪色/\allowbreak{}复色必须解释为旅程或环境变化，不称个人恢复或达标。设备故障由操作员按两组共同程序处理，不必将纯测量故障伪装成A的用户表现。


\reviewsubsection{4.2 原生组的完整体验}
B用场景对象承载实际下发引导G，用独立可分辨的响应表现X，并允许符合规则的累计环境变化。Storm可用目标雨幕风门与实际近景风/\allowbreak{}叶响应；Fade用主潮、补潮、长回流分别表现多步骤，实际补吸只在可识别时出现。原生并不要求隐藏必要信息；形式约束不能妨碍用户分辨提示与自身响应。


基础天气、镜头和环境音保持同源；G变化、X响应与累计变化构成闭环处理。闪电不是成功标记；全局饱和度和累计反馈必须各有明确来源。若Fade保留共同的时间复色，就只作为时间层，B另用局部场景承载数据响应，避免同一像素同时声称时间和表现。


四天气是同一系统的四种应用实例，不各自宣称独立情绪效果。Storm用于普通多步结构，Fade用于漏补吸等复杂事件，先完成这两个切片的设计与真值可行性，再决定四天气完整研究的必要性。


\reviewsubsection{4.3 关键事件的非对称故事板}
SB1 稳定跟随：A按I播放；B在可靠X支持下输出场景响应，G保持或回归I。两组都可形成完整练习，B没有必须持续改节奏的要求。


SB2 被试慢半步：A继续原预设提示，不因X改变；B在确认偏差与合法边界后受限调整下一周期或进入自然呼吸缓冲。实际响应仍显示X，G和I分开记录。不能延长当前停留直到被试追上。


SB3 Fade漏补吸：A仍按固定结构前进；B保留漏步骤事实，不补造Actual、不追补一次强制吸气，在下一合法周期恢复指导。该周期相对I的结构缺口必须计入独立错误统计。


SB4 场景接近模块末尾：A按固定预算结束；B不为等待成功无限延长，不加速剩余动作以补齐次数。记录未回归、少完成的周期及缓冲成本，按共同模块预算结束或预设停止。


SB5 低质量/\allowbreak{}断流：A的预设画面保持开环，测量缺失单独记录；B停止依据无效X调节，使用中性降级路径，暂停数据驱动累计。未知既不等于失败也不等于成功。


SB6 恢复或持续难以跟随：A继续固定方案或按共同停止规则结束；B只在有可用数据与合法边界时恢复适应，达到边界则休息或停止。保留故障与受限事件，不删除这些参与者来制造自适应成功样本。


故事板检验各自方案是否诚实可实现，不再要求A和B提供同一Actual信息或产生相同轨迹。这些是设计规格，尚未成为真实录像证据。


\reviewsection{五、实验、计分与解释}
\reviewsubsection{5.1 两臂设计与流程}
默认参与者间随机分配A/\allowbreak{}B，避免先体验自适应后改变固定组学习。两组佩戴同设备，由相同独立算法用于结果评估；A数据只能进入测量记录，不能反馈到呈现。研究者/\allowbreak{}参与者难以对可见方案盲法，采用中性等预算教学，尽可能让实际呼吸标注和预设评分盲于分组。


流程候选为中性准备与设备校准→PANAS前测→分配揭示与条件教学→相同预算的核心体验→PANAS后测→努力/\allowbreak{}舒适/\allowbreak{}理解及访谈。前测早于条件教学，情绪后测早于阅读考试。教学失败、求助和退出是结果的一部分，不以教到两组同样熟练或排除学不会者来选择样本。


180秒教学与800秒核心仅沿用为工作量候选，须检查自适应缓冲、四结构可识别性与短时量表适用性；不称最佳剂量。没有无项目组，因此两臂比较不能证明整个项目相对休息的绝对效益。


\reviewsubsection{5.2 主要情绪结果与统计候选}
推荐主要候选为PANAS-\allowbreak{}NA后测，以处理前NA校正；估计B减A，负值表示B负性情绪更低。候选模型为NA\_\allowbreak{}post \textasciitilde{} condition + centered\_\allowbreak{}NA\_\allowbreak{}pre + preregistered\_\allowbreak{}strata，以参与者为单位报告估计及双侧95\%区间，HC3方差与缺失方法在新SAP冻结。


正式前须确认中文版本、许可、状态时间指令、分布及最小关注差。如果校园一般样本存在严重地板、测量不能反映目标用途，应在正式研究前修订人群/\allowbreak{}情境或量表，重新论证后冻结；不能看到组间结果后改为理解、正性情绪或最显著的一项。情绪差异跨零不等于等效，不能改用HRV证明情绪获益。


主分析按分配方案包含适应未触发、降级、提前结束者；缺失后测与传感器缺失分开处理。采用事前定义的插补和不利方向敏感性，不默认以前测代填即保守。适应次数、努力、理解、实际暴露等为处理后量，不作为主模型校正变量。


\reviewsubsection{5.3 两套执行参照与共同分母}
对I的结果回答共同练习目标是否达成：按不读取G或分组的实际事件标注，识别步骤顺序、时长及周期。先报告连续时长/\allowbreak{}节奏偏差、缺步骤及可用覆盖，再按预设容差形成理想目标达成候选。匹配允许合法的周期起点平移，避免将相位延迟误作整场失效，但不能用任意时间扭曲把错误对齐成正确。


两组使用相同模块练习预算T。候选覆盖量为该预算中实际完成且符合I的有效周期时间占比；技术信号覆盖与完整结果可判定覆盖分开。T划分为已证实合格C、已证实不合格K与结果未知U，C+K+U=T。短暂断流可能令整个周期未知，不能只把断流秒数当U。报告C/\allowbreak{}T及(C+U)/\allowbreak{}T的操作性界，并列技术/\allowbreak{}结果覆盖与可判定段比例；技术不可观测不标为行为不合格。细则与反例见07候选。跨模块权重在结果前固定，不按B少做后的周期数改分母。


对G的结果回答是否跟得上实际提示：根据真实下发时间线衡量步骤、同步与时序误差，另列为过程结果。G相对I的累计偏离和缓冲占时是控制代价。若对G改善而对I恶化，即使主观轻松，也不能声称共同目标训练性能提高。


现行机会PF以固定Target机会定义，不能原样用于自适应后两组的替代性结论。A-\allowbreak{}02/\allowbreak{}A-\allowbreak{}03/\allowbreak{}S-\allowbreak{}02须重新定义机会、持续时间与未知处理，独立波形标注验证后才能冻结。原0.075非劣界无直接沿用依据；若要提出替代主张，另定有意义的代价界、分析集和足够精度。


\reviewsubsection{5.4 负担、理解与机制检查}
努力、舒适、停止、求助与完整使用成本分别报告，不合成为一个任意总分。PANAS-\allowbreak{}PA、感知控制与投入可作为预先指定的次要或探索结果，避免堆量表稀释主要问题。SCCI只反映原生融入相关操纵，不能检验控制器是否真的使用数据。


旧八题降为形成性蓝图。目标当前/\allowbreak{}下一步骤的共同任务可在语义等值后有限比较；Actual、累计与质量题服务B的可用性验证，A不呈现对应信息时不参与同题正确率排名。对A应验证用户知道画面是预设引导，不把变化当个人成功。撤销上一版12题错误率主终点及其功效演算。


数据驱动性需要独立运行证据：给定相同时钟与环境种子，更换真实输入时A输出不变；B仅在符合适应条件时产生规定差异；低质量时不做有效适应。该检查验证实现，不能证明用户收益或独立机制效应。事后对适应事件分层须承认触发者本来可能更难跟随，不能与未触发者直接作因果比较。


\reviewsubsection{5.5 期待、新奇与实际剂量}
条件说明保持中性，不称B高级或更有效；制作质量与教学预算相当，记录处理前相关经验及预期。体验后的可信度和新奇感只作过程描述，不作主效应校正。两臂整包差异可包含期待与呈现影响，不能据此证明适应机制；单次体验不能证明持续收益。同暴露预算不等于同呼吸练习剂量，B缓冲减少的有效练习时间及结构完成量作为整包代价报告，不能通过校正处理后剂量制造机制结论。

\reviewsubsection{5.6 从统计结果到采用判断}
主情绪结果报告差异及区间，与结果前论证的最小关注差比较；共同I达成、努力、不适及准备成本分别报告。没有事前可接受代价界与足够精度，不称功能非劣或B可替代A，不沿用旧界值。若情绪区间不确定，结论为证据不足；若有相对收益但目标代价或使用成本过大，仍不能直接推荐部署。若情绪无差异而功能更好，只能支持对应功能主张，不更换主要研究成功标准。

\reviewsection{六、科学增量与反证}
\reviewsubsection{6.1 最近工作迫使贡献收敛}
ViBreathe 2021已比较固定、响应式和自适应引导；Gummidela等2021已比较节拍/\allowbreak{}反馈与游戏组合；Yu等CHI 2026已有结合呼吸及HRV的自适应系统和自然形态界面；Chittaro等2024比较真实与替代反馈。因此“首次把数据用于呼吸”“首次自然反馈”及“首次自适应比固定”均不能作为本项目贡献。完整对照及来源见科学附录。


候选增量收敛为：针对多步骤呼吸，在固定机位原生场景中实现受限支持与回归共同理想目标；把情绪、负担、对实际提示的跟随和对共同理想的达成分开，揭示适应带来的收益与目标漂移代价。是否形成可发表新增知识仍取决于更广查新、真实方法质量和结果；换画风或采用新名字不足以成立。


\reviewsubsection{6.2 三项贡献及证据要求}
C1 交互制品与设计方法：给出A/\allowbreak{}B定义、I/\allowbreak{}G/\allowbreak{}X分离、场景状态作用及可解释适应边界。证据为真实数据闭环、规则可检查、失败路径与可重建材料；在实现前称设计候选。若近邻已覆盖该问题，收缩为场景化复现，不能宣称新控制算法。


C2 实证权衡：在共同理想及预算下，报告两完整方案的情绪、执行与负担效应及区间，重点辨别“更容易跟随”是否以放宽目标为代价。证据为随机化、有效测量和独立标注；任何方向均报告，不能按功能或SCCI是否通过隐藏情绪结果。


C3 有限设计知识：基于具体成功/\allowbreak{}失败轨迹和用户解释，说明何时支持、何时回归、何时停止适应，及原生表达的可读边界。独立实例应用支持更广范围；未执行或迁移失败时，只称所测场景的设计经验，不把系统规则自动当普适原则。


\reviewsubsection{6.3 五种必须能诚实解释的结果}
B情绪更好且接近I、负担可接受：支持该整包在目标场景中的相对价值；不能把优势单独归因于原生或自适应。


B更轻松且更能跟G，但明显偏离I：说明支持可能以目标漂移为代价，应限制适应或定位为舒适练习，不能声称既提高训练又改善情绪。


情绪无明确差异：报告区间及精度；独立执行或体验证据按原问题解释，不改主要成功标准。B更差时保留固定方案并定位具体成本。


实际信号不可识别或闭环失效：研究只能评价发生的失败/\allowbreak{}降级方案，不称自适应已实现；优先修复或缩小范围。任何良好动画都不能补出X。


独立应用失败：收缩C3；若最接近研究已覆盖核心增量，则将成果定位为复现与应用验证，重新评估投稿目标，而非修改术语保留创新。


\reviewsection{七、样本、资源与实施衔接}
\reviewsubsection{7.1 样本依据必须跟随新的整包比较}
PANAS作为主要候选可以复用旧量表与SAP思路，但新条件改变方差、效应和缺失，旧功效结果不能直接冻结N。按最小关注原始分差、预试方差/\allowbreak{}前后相关、分配和缺失模拟主要模型；功能代价及边界主张另检查区间精度。不能拿自适应控制器输出的评分作为真值或以100+倒填效果。


仅用于可行性量级：在两臂等分、同一残差标准差$\sigma$、双侧$\alpha=0.05$的正态近似下：
\[
N\approx\frac{4(z_{0.975}+z_{\mathrm{power}})^2}{d_{\mathrm{res}}^2}.
\]
$d_{\mathrm{res}}$为相对于基线调整残差的标准化组间差。$N=100$的95\%区间半宽约$0.392\sigma$，80\%/90\%功效对应$d_{\mathrm{res}}$约0.560/0.648；若正式只剩58--74人，90\%功效约需0.851--0.754。不是预期效应，也不是确认性样本结论。


\reviewsubsection{7.2 人数与工时不重复计算}
操作候选人数账本：12名认知/\allowbreak{}舒适阶段学生，其中8名先做需求访谈，不另加8；真实技术/\allowbreak{}完整流程12–20名不同学生，默认规划20；8位评审中2位尝试独立重建，不另加2。若资源是100名学生，默认正式剩68；若100人还包含8评审，正式剩60。希望正式100学生，默认另需32前置学生与8评审。迁移独立延期。上述不是效度门槛；实际资源含义、前置证据与正式精度决定最终人数，不能倒填效应。


五职责为科学/\allowbreak{}测量、设备/\allowbreak{}状态估计、场景/\allowbreak{}控制策略、用户流程、统计/\allowbreak{}整合，具体人员由团队承接。半个月先评估设备/\allowbreak{}SDK、状态可辨识、双条件切片与真实完整会话成本；没有事实前不承诺产品或正式研究在该期限交付。


\reviewsubsection{7.3 随机化与阶段取舍}
条件隐藏分配仍为两臂；原48块只是两臂24种天气顺序的安排。可变小块及四条Williams天气顺序仍可作为降低执行成本的候选，须根据新流程确认，不因100人自动改X-\allowbreak{}01。固定暴露而非完成格定义研究结束；不按某组成功率补人。


前置阶段分别验证：用户需求与共同目标准确；B可用性及A不误导；真实X与适应边界可辨；完整流程成本与量表分布。Storm/\allowbreak{}Fade切片先检验多步逻辑，四天气是否全部进入正式研究再由资源与覆盖决定，不能自动删掉或增加条件。


\reviewsubsection{7.4 下游先后顺序}
先消费本设计修订R-\allowbreak{}01/\allowbreak{}F-\allowbreak{}02/\allowbreak{}V-\allowbreak{}03与测量；再由P-\allowbreak{}01/\allowbreak{}F-\allowbreak{}05/\allowbreak{}S-\allowbreak{}02/\allowbreak{}U-\allowbreak{}02明确I/\allowbreak{}G/\allowbreak{}X和控制归属，由U12-\allowbreak{}04/\allowbreak{}A-\allowbreak{}03/\allowbreak{}X-\allowbreak{}01/\allowbreak{}U12-\allowbreak{}11形成新协议候选；取得真实修订与冻结后才实施和采集。现有Target不改名复用为两个概念，B自适应不由Unity本地计时器或TD擅自决定。





\reviewsubsection{7.5 本轮操作级设计交付}
\href{05-control-spec.pdf}{05受限自适应规格}给出10条优先级规则、模式、周期边界/即时中断、错相缓冲/速率放慢/回归及参数校准责任；\href{06-paired-storyboards.pdf}{06配对故事板}逐例展开Storm五例、Fade七例，并区分设计图与真实运行证据；\href{07-protocol-candidate.pdf}{07实验候选}补实独立计分、完整未知周期、统一后测、缺失、唯一人数及隐藏8/16联合块候选。

33个直接消费者按operational-design.json交接。新规则、主终点/计分、随机化与阶段人数都仍是下一版协议候选；设计已交付不改变原状态、历史签收或领取身份，数值与正式N必须由真实前置证据确定。

% NORMALIZATION-20261008
\newpage\reviewsection{八、普通学习间隙的规范化边界}
\reviewsubsection{8.1 首个情境与设计状态}
普通校园学习间隙的自愿短时练习；不诱发压力，不按NA筛选。操作员协助，保持两设备与原Unity资产。人群NA较低也保留，不能制造压力来追求显著结果。

18组参数的确定方法、11道理解候选题、两组教学、访谈和完整分析规格已交付于设计索引及05/\allowbreak{}07。唯一20题中文题本与适用使用条件仍未核准；当前不标记设计规范全面闭合，不放行对应真实测量。

\reviewsubsection{8.2 正确的阶段顺序}
设计定稿→对应前置活动批准与材料固定→前置实现及验证→正式参数、分析与协议冻结→正式研究。前置活动有自己的材料/\allowbreak{}程序固定，不等于正式参数先于预试冻结。

\reviewsubsection{8.3 采用、展示与论文}
继续研发：信号及功能可成立、用户任务存在、有可接受修改方向；不要求小样本情绪显著

实践展示：真实A输入隔离与B规则响应/\allowbreak{}失效证据、操作者可执行流程、成本与限制可解释；示意不能当运行演示

论文主张：最近工作差异、预先问题/\allowbreak{}分析、可信真值、正式结果精度和完整失效/\allowbreak{}缺失报告；不保证等级

准备、佩戴、校准、清洁、重置与故障人时单列。两组同佩戴设备的比较不估计现实无设备A相对B的全部净成本。独立迁移继续延期，不承担当前论文必需贡献。


\end{document}
